\chapter{Design UI/UX}
\label{chap:uiux}

\section{Positionnement design du projet}
Le design de la plateforme ne relève pas d'un habillage secondaire ; il participe directement à la lisibilité opérationnelle du système. Dans une application métier de gestion du parc informatique, l'objectif n'est pas seulement de rendre l'interface esthétique, mais de réduire la charge cognitive, de hiérarchiser l'information, de rendre les actions importantes immédiatement visibles et de sécuriser les opérations par la clarté de l'expérience.

Le design retenu adopte un langage visuel premium inspiré des outils SaaS modernes : espaces généreux, cartes structurées, palette contrôlée, contrastes lisibles, typographie hiérarchisée, icônes cohérentes et tableaux simples à parcourir. Cette direction renforce la perception de qualité du produit et facilite la démonstration académique.

\section{Design system}
Le \textit{design system} de la plateforme s'appuie sur plusieurs principes :
\begin{itemize}
  \item uniformité des composants interactifs ;
  \item hiérarchisation visuelle stable des titres, sous-titres et textes de support ;
  \item usage récurrent de cartes, badges, tableaux et filtres ;
  \item feedback visuel cohérent pour les états, les alertes et les statuts métier.
\end{itemize}

Ce système évite l'effet de dispersion graphique. Il rend l'application plus professionnelle et plus simple à faire évoluer.

\section{Palette de couleurs}
La palette graphique privilégie des tons institutionnels et professionnels :
\begin{itemize}
  \item un rouge bordeaux pour les éléments identitaires et certains accents forts ;
  \item des bleus doux pour les surfaces d'information et les états neutres ;
  \item des verts et ambres pour les statuts positifs, préventifs ou en surveillance ;
  \item des gris très clairs pour le fond général et la respiration visuelle.
\end{itemize}

Ce choix est pertinent pour une application de gestion, car il permet de conserver un ton sérieux tout en évitant la monotonie des interfaces purement techniques.

\section{Typographie}
La plateforme mobilise une hiérarchie typographique orientée lisibilité :
\begin{itemize}
  \item une police d'interface claire pour le corps du texte ;
  \item une police à forte présence pour les titres et accroches ;
  \item des niveaux de graisse différenciés pour distinguer les métadonnées, les labels et les indicateurs clés.
\end{itemize}

Cette hiérarchie est essentielle dans les tableaux de bord, où l'utilisateur doit distinguer rapidement signal faible, indicateur principal et action disponible.

\section{Expérience utilisateur}
L'expérience utilisateur a été pensée selon plusieurs axes :
\begin{enumerate}
  \item \textbf{accès rapide} : la page de connexion met immédiatement en valeur l'objectif du produit tout en réduisant la friction ;
  \item \textbf{navigation stable} : le menu latéral hiérarchise les modules par familles d'usage ;
  \item \textbf{visualisation synthétique} : le dashboard présente les informations clés avant les détails ;
  \item \textbf{progressivité} : l'utilisateur passe d'une vue synthétique à une vue opérationnelle sans rupture graphique.
\end{enumerate}

Dans une logique UI/UX, ce type de structure est préférable à une multiplication de pages hétérogènes sans fil conducteur.

Une amélioration récente concerne l'espace employé bénéficiaire: la page des affectations devient également un espace \enquote{Mes matériels}. L'utilisateur peut y consulter ses dotations, déposer une demande de nouveau matériel et suivre le statut de ses demandes. Cette logique rend le parcours plus naturel, car le bénéficiaire n'a plus besoin de comprendre les mecanismes administratifs d'affectation pour exprimer son besoin.

\section{Responsive design}
Le responsive design a été intégré comme exigence réelle, non comme simple adaptation tardive. Les captures générées montrent que l'application reste exploitable sur tablette et sur des largeurs plus contenues. Les composants se réorganisent, les cartes se superposent verticalement et la navigation conserve son rôle principal. Les tableaux et le Kanban de maintenance ont aussi été ajustés pour contenir leur défilement horizontal à l'intérieur des composants, ce qui évite les décalages de page constatés lors des tests visuels.

\section{Corrections ergonomiques finales}
Une dernière passe de stabilisation UI a été réalisée après les tests utilisateurs. Elle a porte sur les points suivants:
\begin{itemize}
  \item mise en avant du formulaire \enquote{Demander un matériel} dans l'espace employé bénéficiaire, afin que l'action principale soit visible sans recherche;
  \item transformation des grilles trop serrees en mises en page empilees sur les écrans moyens, notamment pour les modules équipements, utilisateurs et maintenance;
  \item ajout d'un défilement horizontal interne aux tableaux pour éviter que la page complète ne se décale vers la droite;
  \item reduction de là largeur minimale du Kanban et stabilisation des colonnes pour supprimer l'effet d'écran noir ou de deplacement du contenu;
  \item reorganisations des formulaires longs en grilles lisibles, avec une colonne mobile et deux colonnes desktop lorsque l'espace le permet.
\end{itemize}

Ces corrections renforcent l'objectif principal du design: conserver une interface claire, premium et utilisable dans des conditions réelles de démonstration, même lorsque les tableaux contiennent beaucoup de colonnes ou que l'écran est plus etroit.

\section{Accèssibilité}
Même si le projet n'implémente pas un audit d'accessibilité complet, plusieurs bonnes pratiques sont visibles :
\begin{itemize}
  \item labels explicites sur les champs de formulaire ;
  \item contrastes suffisants sur les éléments importants ;
  \item tailles de titres et zones interactives confortables ;
  \item structuration logique des sections.
\end{itemize}

Ces choix contribuent à une meilleure robustesse d'usage et facilitent la prise en main dans un contexte professionnel varié.

\section{Captures d'écran principales}
\subsection{Dashboard}
Le dashboard représente la vue centrale de pilotage. Il combine indicateurs, graphiques, activité récente et résumé de maintenance.

\screenfigure{dashboard.png}{Dashboard principal}{fig:dashboard}

Le dashboard met en oeuvre une logique de lecture en Z : résumé chiffré en haut, visualisation au centre, détails opérationnels en bas. Cette construction améliore la perception immédiate de l'état du parc.

\subsection{Gestion des équipements}
La gestion des équipements est l'un des modules les plus sensibles, car elle conditionne la qualité de toutes les autres opérations.

\screenfigure{equipments.png}{Module de gestion des équipements}{fig:equipments}

L'écran privilégie la lisibilité tabulaire, la hiérarchisation des statuts et l'accès rapide aux actions. Cette orientation est adaptée à un contexte d'inventaire, où la densité d'information doit rester maîtrisée.

\subsection{Affectations et utilisateurs}
Les modules utilisateurs et affectations jouent un rôle majeur dans la gouvernance du parc.

\screenfigure{assignments.png}{Module des affectations}{fig:assignments}

La cohérence graphique avec la gestion des utilisateurs renforce la sensation d'un produit unifié. Les modèles de tableaux, filtres, badges et en-têtes sont réutilisés intelligemment. La page des affectations intégré desormais une zone de demandes de matériel: formulaire côté employé et file de traitement côté administration.

\subsection{Maintenance, tickets et Kanban}
La dimension maintenance bénéficie d'une visualisation spécifique, adaptée aux flux de traitement.

\screenfigure{maintenance.png}{Vue maintenance}{fig:maintenance}

La représentation par état, inspirée des tableaux Kanban, apporte un bénéfice concret : elle rend visible la charge et facilite la priorisation des actions. Une correction de mise en page a été appliquée pour stabiliser les colonnes, garder une largeur lisible pour chaque carte et éviter le décalage du contenu principal.

\subsection{Notifications et rapports}
Les modules notifications et rapports renforcent la dimension de pilotage et de communication.
Ils sont conservés dans le périmètre fonctionnel du projet, mais leurs captures détaillées ne sont pas reproduites afin de limiter le poids du PDF final.

\section{Responsive mobile et tablette}
Le comportement responsive peut être observé dans les figures \ref{fig:mobile-dashboard-1} et \ref{fig:mobile-equipments-1}. L'adaptation mobile ne se limite pas à une réduction de taille ; elle réorganise réellement les zones fonctionnelles.

\mobilefigure{dashboard-mobile-part-1.png}{Dashboard mobile : en-tête et indicateurs}{fig:mobile-dashboard-1}

\mobilefigure{equipments-mobile-part-1.png}{Gestion des équipements mobile : liste et filtres}{fig:mobile-equipments-1}

\section{Évaluation critique du design}
Sur le plan critique, plusieurs forces ressortent clairement :
\begin{itemize}
  \item l'identité visuelle est homogène ;
  \item les modules principaux sont immédiatement identifiables ;
  \item la plateforme donne une impression de produit réellement exploitable ;
  \item la qualité perçue dépasse celle d'un prototype académique basique.
\end{itemize}

Quelques pistes d'amélioration peuvent néanmoins être envisagées :
\begin{itemize}
  \item enrichir les interactions de confirmation sur certaines actions critiques ;
  \item approfondir l'accessibilité clavier et lecteur d'écran ;
  \item proposer davantage de personnalisation sur les tableaux et filtres ;
  \item intégrer des modes de visualisation encore plus avancés pour les analytics.
\end{itemize}

\section{Synthèse du chapitre}
Le design UI/UX de la plateforme contribue directement à sa valeur. Il ne s'agit pas d'un habillage cosmétique, mais d'un levier d'adoption, de lisibilité métier et de crédibilité professionnelle. Les captures d'écran réelles montrent une interface moderne, cohérente et bien hiérarchisée, compatible avec les objectifs académiques et opérationnels du projet.









